%---------------------------preamble--------------------------------------
%-------------------------------------------------------------------------
%--------------Preamble----------------------------------------------------
%compile this package with pdflatex
\documentclass[10pt]{beamer} % default
%\documentclass[10pt,aspectratio=169]{beamer} % default
%\documentclass[draft,10pt,aspectratio=169]{beamer} % draft
%\documentclass[handout,10pt,aspectratio=169]{beamer} % handout
%-------------------------------------------------------------------------
% my def
\definecolor{tuhh_darkcyan}{rgb}{.09,0.39,.42}
\def\LBOX{\fbox{\begin{minipage}[c][.4cm][c]{.5cm}%
  \center \large \phantom{1} \qquad\end{minipage}}}
\newcommand\BOXWF[1]{\fbox{\begin{minipage}[c][.4cm][c]{.5cm}%
  \center \large {#1}\end{minipage}}}
% my def
%-------------------------------------------------------------------------
%
\usepackage{amsfonts}
\usepackage{amssymb}
\usepackage[ngerman]{babel}
\usepackage{ulem}
\normalem
\usepackage{bookmark}
\usepackage{apalike}
%%%%%%%%%%%%%%%%%%%%%%%%%%%%%%%%%%%%%%%%%%%%%%%%%%%%%%%%%%%%%%%%%%%%%%%%%%%%%%%
% Definitions
%%%%%%%%%%%%%%%%%%%%%%%%%%%%%%%%%%%%%%%%%%%%%%%%%%%%%%%%%%%%%%%%%%%%%%%%%%%%%%%
%
% Colors:
  \definecolor{myalgcolor}{rgb}{1.0,1.0,0.5}
  \definecolor{myred}{rgb}{1.0,0.0,0.0}
  \definecolor{myyellow}{rgb}{1.0,1.0,0.5}
  \definecolor{mygray}{gray}{0.55}
  \definecolor{mycyan}{rgb}{.0,.8,1.0}
  \definecolor{blau}{rgb}{.5,.5,1.0}
  \definecolor{myerror}{rgb}{1.0,0.0,0.5}
  \definecolor{myerron}{rgb}{0.4,0.0,0.6}
\newcommand{\error}[1]{\textcolor{myerror}{{}#1{}}}
\newcommand{\erron}[1]{\textcolor{myerron}{{}#1{}}}
\newcommand{\backupbegin}{
   \newcounter{framenumberappendix}
   \setcounter{framenumberappendix}{\value{framenumber}}
}
\newcommand{\backupend}{
   \addtocounter{framenumberappendix}{-\value{framenumber}}
   \addtocounter{framenumber}{\value{framenumberappendix}} 
}
\DeclareMathOperator{\diag}{\mathsf{diag}}
\DeclareMathOperator{\argmin}{\mathsf{argmin}}
\DeclareMathOperator{\codim}{\mathsf{codim}}
\DeclareMathOperator{\Rang}{\mathsf{Rang}}
\DeclareMathOperator{\zeros}{\mathsf{zeros}}
\newcommand{\bfnu}{\boldsymbol{\nu}}
\newcommand{\bfhatnu}{\boldsymbol{\widehat{\nu}}}
\newcommand{\bfchecknu}{\boldsymbol{\check{\nu}}}
\DeclareMathOperator{\sut}{\mathsf{sut}}
\DeclareMathOperator{\tril}{\mathsf{tril}}
\DeclareMathOperator{\triu}{\mathsf{triu}}
\renewcommand{\diag}{\mathsf{diag}}
\DeclareMathOperator{\diagm}{\mathsf{diag}}
\DeclareMathOperator{\trace}{\mathsf{trace}}
\DeclareMathOperator{\sign}{\mathsf{sign}}
%
%%%%%%%%%%%%%%%%%%%%%%%%%%%%%%%%%%%%%%%%%%%%%%%%%%%%%%%%%%%%%%%%%%%%%%%%%%%%%%%
% Definitions added by Jens-Peter Zemke
%%%%%%%%%%%%%%%%%%%%%%%%%%%%%%%%%%%%%%%%%%%%%%%%%%%%%%%%%%%%%%%%%%%%%%%%%%%%%%%
%% my own stuff
%%
%% definition of perlis mathclap etcpp.
%%
% For comparison, the existing overlap macros:
% \def\llap#1{\hbox to 0pt{\hss#1}}
% \def\rlap#1{\hbox to 0pt{#1\hss}}
\def\clap#1{\hbox to 0pt{\hss#1\hss}}
\def\mathllap{\mathpalette\mathllapinternal}
\def\mathrlap{\mathpalette\mathrlapinternal}
\def\mathclap{\mathpalette\mathclapinternal}
\def\mathllapinternal#1#2{%
  \llap{$\mathsurround=0pt#1{#2}$}}
\def\mathrlapinternal#1#2{%
  \rlap{$\mathsurround=0pt#1{#2}$}}
\def\mathclapinternal#1#2{%
  \clap{$\mathsurround=0pt#1{#2}$}}
%%
%% end of perlis stuff
%%
%%
%% definition of inverse diagonal dots (iddots)
%%
\makeatletter
\def\iddots{\mathinner{\mkern1mu\raise\p@%
    \hbox{.}\mkern2mu\raise4\p@\hbox{.}\mkern2mu%
    \raise7\p@\vbox{\kern7\p@\hbox{.}}\mkern1mu}}
\makeatother
%%%\def\ddots{\mathinner{\mkern1mu\raise7\p@\vbox%
%%%{\kern7\p@\hbox{.}} \mkern2mu\raise4\p@\hbox{.}%
%%%\mkern2mu\raise\p@\hbox{.}\mkern1mu}}
%
% TIME OF DAY (Nelson H. F. Beebe :-)
%
\newcount\hh
\newcount\mm
\mm=\time
\hh=\time
\divide\hh by 60
\divide\mm by 60
\multiply\mm by 60
\mm=-\mm
\advance\mm by \time
\def\hhmm{\number\hh:\ifnum\mm<10{}0\fi\number\mm}
%
% Use it like this in a LaTeX document:
%
%        \date{\today{ }\hhmm}
%
%%%%%%%%%%%%%%%%%%%%%%%%%%%%%%%%%%%%%%%%%%%%%%%%%%%%%%%%%%%%%%%
%%%
%%% kleiner Hack, um Einträge in Matrizen anders auszurichten
%%%
%%%%%%%%%%%%%%%%%%%%%%%%%%%%%%%%%%%%%%%%%%%%%%%%%%%%%%%%%%%%%%%
\makeatletter
\renewcommand*\env@matrix[1][*\c@MaxMatrixCols c]{%
  \hskip -\arraycolsep
  \let\@ifnextchar\new@ifnextchar
  \array{#1}}
\makeatother
%%%%%%%%%%%%%%%%%%%%%%%%%%%%%%%%%%%%%%%%%%%%%%%%%%%%%%%%%%%%%%%
%%%
%%% Wir wollen Matlab-Beispielprogramme einbinden
%%%
%%%%%%%%%%%%%%%%%%%%%%%%%%%%%%%%%%%%%%%%%%%%%%%%%%%%%%%%%%%%%%%
\usepackage{listings}
\definecolor{hellgrau}{rgb}{0.90,0.90,0.90}
\definecolor{commentcol}{rgb}{0.0823,.4902,0.0}
\lstset{language=Matlab,
        basicstyle={\footnotesize\ttfamily},
        keywordstyle={\sffamily\bfseries},
        tabsize=2,
        escapechar=\#,
        numbers=left,
        numberstyle=\tt,
        stepnumber=1,
        numbersep=7pt,
        breaklines=true,
        frame=single,
        frameround=ffff,
        commentstyle=\color{commentcol},
        backgroundcolor=\color{hellgrau}
}
\usepackage{algorithmic}
\renewcommand{\algorithmiccomment}[2]{\hfill\rlap{\texttt{\%
      #1}}\phantom{\texttt{\% #2}}}
\renewcommand{\algorithmicrequire}{\textsc{Eingabe:}}
\renewcommand{\algorithmicensure}{\textsc{Ausgabe:}}
\usepackage{algorithm}
\floatname{algorithm}{Algorithmus}

\mode<presentation>
{
  \usetheme{ABMATH}
  \useinnertheme{default}
  \setbeamercovered{invisible}
  \usefonttheme[onlymath]{serif}
  \setbeamertemplate{items}[triangle]
  \setbeamertemplate{sections/subsections in toc}[default]
  \setbeamercolor*{author in head/foot}{parent=palette secondary}
}
\setbeamerfont{quote}{shape=\upshape}
\usepackage{pgf}
\usepackage[utf8]{inputenc}
\usepackage[T1]{fontenc}
\usepackage{times}
\usepackage{graphicx}
\usepackage{pifont}
\usepackage{colortbl}
\usepackage{xmpmulti}
\usepackage{multimedia}
\usepackage{pdfpages}
\usepackage{tikz}
\usepackage{pgfplots}

\newcommand{\bfA}{{\mathbf A}}
\newcommand{\bfB}{{\mathbf B}}
\newcommand{\bfC}{{\mathbf C}}
\newcommand{\bfD}{{\mathbf D}}
\newcommand{\bfE}{{\mathbf E}}
\newcommand{\bfF}{{\mathbf F}}
\newcommand{\bfG}{{\mathbf G}}
\newcommand{\bfH}{{\mathbf H}}
\newcommand{\bfI}{{\mathbf I}}
\newcommand{\bfJ}{{\mathbf J}}
\newcommand{\bfK}{{\mathbf K}}
\newcommand{\bfL}{{\mathbf L}}
\newcommand{\bfM}{{\mathbf M}}
\newcommand{\bfN}{{\mathbf N}}
\newcommand{\bfO}{{\mathbf O}}
\newcommand{\bfP}{{\mathbf P}}
\newcommand{\bfQ}{{\mathbf Q}}
\newcommand{\bfR}{{\mathbf R}}
\newcommand{\bfS}{{\mathbf S}}
\newcommand{\bfT}{{\mathbf T}}
\newcommand{\bfU}{{\mathbf U}}
\newcommand{\bfV}{{\mathbf V}}
\newcommand{\bfW}{{\mathbf W}}
\newcommand{\bfX}{{\mathbf X}}
\newcommand{\bfY}{{\mathbf Y}}
\newcommand{\bfZ}{{\mathbf Z}}
\newcommand{\bfa}{{\mathbf a}}
\newcommand{\bfb}{{\mathbf b}}
\newcommand{\bfc}{{\mathbf c}}
\newcommand{\bfd}{{\mathbf d}}
\newcommand{\bfe}{{\mathbf e}}
\newcommand{\bff}{{\mathbf f}}
\newcommand{\bfg}{{\mathbf g}}
\newcommand{\bfh}{{\mathbf h}}
\newcommand{\bfi}{{\mathbf i}}
\newcommand{\bfj}{{\mathbf j}}
\newcommand{\bfk}{{\mathbf k}}
\newcommand{\bfl}{{\mathbf l}}
\newcommand{\bfm}{{\mathbf m}}
\newcommand{\bfn}{{\mathbf n}}
\newcommand{\bfo}{{\mathbf o}}
\newcommand{\bfp}{{\mathbf p}}
\newcommand{\bfq}{{\mathbf q}}
\newcommand{\bfr}{{\mathbf r}}
\newcommand{\bfs}{{\mathbf s}}
\newcommand{\bft}{{\mathbf t}}
\newcommand{\bfu}{{\mathbf u}}
\newcommand{\bfv}{{\mathbf v}}
\newcommand{\bfw}{{\mathbf w}}
\newcommand{\bfx}{{\mathbf x}}
\newcommand{\bfy}{{\mathbf y}}
\newcommand{\bfz}{{\mathbf z}}
\newcommand{\bfell}{{\boldsymbol\ell}}
\newcommand{\bfSigma}{{\boldsymbol\Sigma}}
\newcommand{\DEF}{:=}
\newcommand{\FED}{=:}
\newcommand{\DEFM}[1]{\textcolor{red}{#1}}
\newcommand{\T}{\mathsf{T}}
\renewcommand{\H}{\mathsf{H}}
\newcommand{\Span}{\mathop{\mathsf{span}}}
\newcommand{\spec}{\mathop{\mathsf{spec}}}

%%% thanks, Nick Higham :)
\setbeamertemplate{navigation symbols}{}

\usepackage[absolute,overlay]{textpos}
%%%%%%%%%%%%%%%%%%%%%%%%%%%%%%%%%%%%%%%%%%%%%%%%%%%%%%%%%%%%%%
%%%%%%%%%%%% my own packages %%%%%%%%%%%%%%%%%%%%%%%%%%%%%%%%%
%%%%%%%%%%%%%%%%%%%%%%%%%%%%%%%%%%%%%%%%%%%%%%%%%%%%%%%%%%%%%%

\title[Kurzer Titel]%
{Vom Sinn und Nutzen eines Proseminares,
  auch bekannt als <langer Titel>}

\author[Jens-Peter M. Zemke]%
{%
  \href{https://www.mat.tuhh.de/home/jpmzemke}%
          {Jens-Peter M. Zemke}\break
  \footnotesize{\href{mailto:zemke@tu-harburg.de}%
                {zemke@tu-harburg.de}}%
}

\titlegraphic{%
  \includegraphics[height=.15\textheight]%
    {pics/TUHH_de}%
}

\institute{\href{https://www.mat.tuhh.de/}{%
            Institut für Mathematik}\\
            Lehrstuhl Numerische Mathematik\\
           \href{https://www.tuhh.de/}{%
            Technische Universität Hamburg}%
}

\date[2017-04-03]{\footnotesize{3. April 2017}}

\AtBeginSubsection[]
{
  \begin{frame}<beamer>
    \frametitle{Übersicht}
    \tableofcontents[currentsection,currentsubsection]
%    \tableofcontents[currentsection]
  \end{frame}
}
\AtBeginSection[]
{
  \begin{frame}<beamer>
    \frametitle{Übersicht}
%    \tableofcontents[currentsection,currentsubsection]
    \tableofcontents[currentsection]
  \end{frame}
}
%---------------------------preamble-----------------------------------
\hypersetup{%
  pdftitle     = {},
  pdfsubject   = {},
  pdfkeywords  = {},
  pdfauthor    = {\textcopyright\ Jens-Peter M. Zemke, März 2017}, 
  % pdfpagemode  = FullScreen
%  colorlinks = true,
%  linkcolor = ,
%  urlcolor = blue,
}
%%%%%%%%%%%%%%%%%%%%%%%%%%%%%%%%%%%%%%%%%%%%%%%%%%%%%%%%%%%%%%%%%%%%%%%%%%%%%%%
% Frontmatter
%%%%%%%%%%%%%%%%%%%%%%%%%%%%%%%%%%%%%%%%%%%%%%%%%%%%%%%%%%%%%%%%%%%%%%%%%%%%%%%

\begin{document}

%----------------------------begin frame-------------------------------
\begin{frame}
  \titlepage
\end{frame}
%----------------------------end frame---------------------------------
\usebackgroundtemplate{%
      \includegraphics[width=\paperwidth]{pics/Background}}
%----------------------------begin frame------------------------------
\begin{frame}
  \frametitle{Übersicht}
  \tableofcontents[pausesections]
\end{frame}
%----------------------------end frame--------------------------------
\section{Grundlegendes}
%----------------------------begin frame------------------------------
\begin{frame}
  \frametitle{Das Wichtigste in Kürze}

  Eckdaten:
  \begin{itemize}
  \item<1-> Titel: „Proseminar über numerische lineare Algebra“
  \item<2-> Literatur: Buch \cite{Watkins:2002} von David Watkins
  \item<3-> Teilnehmer: maximal 36 Abschnitte $\leadsto$ maximal 36
  \item<4-> Treffen: Montag, 17:00--18:30, Channel Harburg, CH4-0.08
  \item<5-> 3 Vorträge pro Treffen
  \end{itemize}

  \vspace*{1em}

\uncover<6->{%
  Anforderungen:
  \begin{itemize}
  \item<7-> \LaTeX-Ausarbeitung (3--15 Seiten)
  \item<8-> \textsc{Matlab}-Beispielprogramme
  \item<9->  Vortrag auf Basis der \LaTeX-Beamer-Klasse (30\,min)
  \item<10-> Anwesenheit bei den Vorträgen
  \end{itemize}}

\end{frame}
%----------------------------end frame--------------------------------
\subsection{Das mathematische Thema}
%----------------------------begin frame------------------------------
\begin{frame}
  \frametitle{Thema}
  
  Thema $\leadsto$ erste Kapitel des Buches~\cite{Watkins:2002},
  “Fundamentals of Matrix Computations”:

  \vspace*{1em}

  \begin{itemize}[<+->]
  \item Gauß'sches Eliminationsverfahren,
  \item Sensitivität linearer Gleichungssysteme,
  \item Ausgleichsprobleme,
  \item Singulärwertzerlegung (SVD),
  \item Eigenwerten und Eigenvektoren,
  \item (iterative Methoden zur Lösung linearer Gleichungssysteme).
  \end{itemize}

\end{frame}
%----------------------------end frame--------------------------------
%----------------------------begin frame------------------------------
\begin{frame}
  \frametitle{Kapitel 1: Gauß/LR}

  \begin{tabular}{|r|l|l|}\hline
  \textbf{No.} & \textbf{Name} & \textbf{Thema} \\\hline
  \multicolumn{3}{|l|}{Kapitel 1: Gaussian Elimination and Its Variants.}\\\hline
    \textcolor{blue}{1}&&
  1.1 Matrix Multiplication.\\\hline
    \textcolor{blue}{2}&&
  1.2 Systems of Linear Equations.\\\hline
    \textcolor{blue}{3}&&
  1.3 Triangular Systems.\\\hline
    \textcolor{red}{4}&&
  1.4 Positive Definite Systems; Cholesky Decomposition.\\\hline
    \textcolor{red}{5}&&
  1.5 Banded Positive Definite Systems.\\\hline
    \textcolor{red}{6}&&
  1.6 Sparse Positive Definite Systems.\\\hline
    \textcolor{black}{7}&&
  1.7 Gaussian Elimination and the LU Decomposition.\\\hline
    \textcolor{black}{8}&&
  1.8 Gaussian Elimination with Pivoting.\\\hline
    \textcolor{black}{9}&&
  1.9 Sparse Gaussian Elimination.\\\hline
  \end{tabular}

\end{frame}
% ----------------------------end frame--------------------------------
%----------------------------begin frame------------------------------
\begin{frame}
  \frametitle{Kapitel 2: Sensitivität von $\bfA\bfx=\bfb$}

  \begin{tabular}{|r|l|l|}\hline
  \textbf{No.} & \textbf{Name} & \textbf{Thema} \\\hline
  \multicolumn{3}{|l|}{Kapitel 2: Sensitivity of Linear Systems.}\\\hline
    \textcolor{blue}{10}&&
  2.1 Vector and Matrix Norms.\\\hline
    \textcolor{blue}{11}&&
  2.2 Condition Numbers.\\\hline
    \textcolor{blue}{12}&&
  2.3 Perturbing the Coefficient Matrix.\\\hline
    \textcolor{red}{13}&&
  2.4 A Posteriori Error Analysis Using the Residual.\\\hline
    \textcolor{red}{14}&&
  2.5 Roundoff Errors; Backward Stability.\\\hline
    \textcolor{red}{15}&&
  2.6 Propagation of Roundoff Errors.\\\hline
    \textcolor{black}{16}&&
  2.7 Backward Error Analysis of Gaussian Elimination.\\\hline
    \textcolor{black}{17}&&
  2.8 Scaling.\\\hline
    \textcolor{black}{18}&&
  2.9 Componentwise Sensitivity Analysis.\\\hline
  \end{tabular}

\end{frame}
% ----------------------------end frame--------------------------------
%----------------------------begin frame------------------------------
\begin{frame}
  \frametitle{Kapitel 3: Ausgleichsprobleme $\min\|\bfA\bfx-\bfb\|_2$}
  
  \begin{tabular}{|r|l|l|}\hline
  \textbf{No.} & \textbf{Name} & \textbf{Thema} \\\hline
  \multicolumn{3}{|l|}{Kapitel 3: The Least Squares Problem.}\\\hline
    \textcolor{blue}{19}&&
  3.1 The Discrete Least Squares Problem.\\\hline
    \textcolor{blue}{20}&&
  3.2 Orthogonal Matrices, Rotators and Reflectors.\\\hline
    \textcolor{blue}{21}&&
  3.3 Solution of the Least Squares Problem.\\\hline
    \textcolor{red}{22}&&
  3.4 The Gram-Schmidt Process.\\\hline
    \textcolor{red}{23}&&
  3.5 Geometric Approach.\\\hline
    \textcolor{red}{24}&&
  3.6 Updating the QR Decomposition.\\\hline
  \end{tabular}

\end{frame}
% ----------------------------end frame--------------------------------
%----------------------------begin frame------------------------------
\begin{frame}
  \frametitle{Kapitel 4: SVD $\bfA=\bfU\bfSigma\bfV^{\mathsf{H}}$}
  
  \begin{tabular}{|r|l|l|}\hline
  \textbf{No.} & \textbf{Name} & \textbf{Thema} \\\hline
  \multicolumn{3}{|l|}{Kapitel 4: The Singular Value Decomposition.}\\\hline
    \textcolor{black}{25}&&
  4.1 Introduction.\\\hline
    \textcolor{black}{26}&&
  4.2 Some Basic Applications of Singular Values.\\\hline
    \textcolor{black}{27}&&
  4.3 The SVD and the Least Squares Problem.\\\hline
    \textcolor{blue}{28}&&
  4.4 Sensitivity of the Least Squares Problem.\\\hline
  \end{tabular}

\end{frame}
% ----------------------------end frame--------------------------------
%----------------------------begin frame------------------------------
\begin{frame}
  \frametitle{Kapitel 5: Eigenwerte und Eigenvektoren
    $\bfA\bfv=\bfv\lambda$}

  \begin{tabular}{|r|l|l|}\hline
  \textbf{No.} & \textbf{Name} & \textbf{Thema} \\\hline
  \multicolumn{3}{|l|}{Kapitel 5: Eigenvalues and Eigenvectors I.}\\\hline
    \textcolor{blue}{29}&&
  5.1 Systems of Differential Equations.\\\hline
    \textcolor{blue}{30}&&
  5.2 Basic Facts.\\\hline
    \textcolor{red}{31}&&
  5.3 The Power Method and Some Simple Extensions.\\\hline
    \textcolor{red}{32}&&
  5.4 Similarity Transforms.\\\hline
    \textcolor{red}{33}&&
  5.5 Reduction to Hessenberg and Tridiagonal Forms.\\\hline
    \textcolor{black}{34}&&
  5.6 The QR Algorithm.\\\hline
    \textcolor{black}{35}&&
  5.7 Use of the QR Algorithm to Calculate Eigenvectors.\\\hline
    \textcolor{black}{36}&&
  5.8 The SVD Revisted.\\\hline
  \end{tabular}

\end{frame}
%----------------------------end frame--------------------------------
\subsection{Die Ausarbeitung in \LaTeX}
%----------------------------begin frame------------------------------
\begin{frame}
  \frametitle{\TeX, \LaTeX, Formeln}

  \begin{description}
  \item[\LaTeX:] Makros von Leslie Lamport zu
  \item[\TeX:] Textsatzsystem von Donald Knuth
  \end{description}

  \vspace*{1em}

\uncover<2->{%
  \LaTeX\ perfekt geeignet für mathematische Texte:
  \begin{itemize}
  \item Formeln: $f(x):=x^2-\sin(x)$,
  \item Terme: $\int_0^1f(x)\,dx$.
  \end{itemize}}

  \vspace*{1em}

\uncover<3->{%
  Abgesetzte nummerierte Formeln:
  \begin{equation}
    \label{eq:eine_erste_abgesetzte_Formel}
    \begin{pmatrix}
      1 & 2 & 3\\
      4 & 5 & 6\\
      7 & 8 & 9\\
    \end{pmatrix}
    \begin{pmatrix}[r]
      1 \\ -2 \\ 1\\
    \end{pmatrix} = 
    \begin{pmatrix}
      0 \\ 0 \\ 0\\
    \end{pmatrix} \quad\Leftrightarrow\quad
    \bfA\bfx=\bfo.
  \end{equation}}

\uncover<4->{%
  \begin{itemize}
  \item Zugriff mittels \texttt{\textbackslash eqref}:
    Formel~\eqref{eq:eine_erste_abgesetzte_Formel},
  \item schöner mittels Hinzunahme von \texttt{\textbackslash
      hyperref}:
         \hyperref[eq:eine_erste_abgesetzte_Formel]%
    {Formel~(\ref*{eq:eine_erste_abgesetzte_Formel})}.
  \end{itemize}}

\end{frame}
% ----------------------------end frame--------------------------------
%----------------------------begin frame------------------------------
\begin{frame}
  \frametitle{Formeln \& Nummerierung}

  Mehrere Formeln gruppieren $\leadsto$ \texttt{align}:
  \begin{align}
    \label{eq:erste_Zeile_gruppierter_Formeln}
    a(x) &= x^2+4, & b(x) &= x^3-4x^2, \\
    \label{eq:zweite_Zeile_gruppierter_Formeln}
    c(x) &= x^2+x, & d(x) &= x^3+13, \\
    \label{eq:dritte_Zeile_gruppierter_Formeln}
    e(x) &= x^2+4x, & f(x) &= x^5-12x^3.
  \end{align}
\uncover<2->{%
  Alle Zeilen eigene Formelnummer, unterdrücken mittels
  \texttt{\textbackslash notag}.

  \begin{enumerate}
  \item die erste Zeile hat die
    Nummer~\eqref{eq:erste_Zeile_gruppierter_Formeln},
  \item die zweite Zeile hat die
    Nummer~\eqref{eq:zweite_Zeile_gruppierter_Formeln}.
  \end{enumerate}}

\uncover<3->{%
  Nummern zusammenfassen $\leadsto$ \texttt{subequations}:
  \begin{subequations}
  \label{eq:gruppierte_Formeln}
    \begin{align}
      \label{eq:erste_Zeile_gruppierter_Formeln_a}
      a(x) &= x^2+4, & b(x) &= x^3-4x^2, \\
      \label{eq:zweite_Zeile_gruppierter_Formeln_b}
      c(x) &= x^2+x, & d(x) &= x^3+13, \\
      \label{eq:dritte_Zeile_gruppierter_Formeln_c}
      e(x) &= x^2+4x, & f(x) &= x^5-12x^3.
    \end{align}
  \end{subequations}

  Zweite Zeile:~\eqref{eq:zweite_Zeile_gruppierter_Formeln_b},
  gesamter Formelblock:~\eqref{eq:gruppierte_Formeln}.}

\end{frame}
%----------------------------end frame--------------------------------
%----------------------------begin frame------------------------------
\begin{frame}
  \frametitle{Tabellen, mehrspaltiges Layout}

  \begin{columns}[onlytextwidth]
    \column{.5\textwidth}
    Tabellen mittels \texttt{tabular}, wie
    \begin{center}
      \begin{tabular}{|l|c|}\hline
        \textbf{Name} & \textbf{Gruppe} \\\hline\hline
        Eduard & 4 \\\hline
        Alexander & 2 \\\hline
      \end{tabular},
    \end{center}
    sind auch möglich.

   \vspace*{1em}

\uncover<2->{%
    Mathematische Tabellen sind mittels \texttt{array} im Mathemodus
    zu setzen,
    \[
      \begin{array}{|c|l|}\hline
        n/2 & \text{$n$ gerade}\\\hline
        3n+1 & \text{$n$ ungerade}\\\hline
      \end{array}.
    \]}

    \column{.5\textwidth}
\uncover<3->{%
    Es ist ziemlich leicht, in Beamer mehrspaltigen Text zu
    setzen. Dieser Text ist in der zweiten Spalte.}
  \end{columns}

  \vspace*{1em}

\uncover<4->{%
  Dieser (etwas längere) Text ist außerhalb des zweispaltigen
  Textsatzes.}

\end{frame}
% ----------------------------end frame--------------------------------
%----------------------------begin frame------------------------------
\begin{frame}[fragile]
  \frametitle{\LaTeX\ \& Bilder}

  In \LaTeX\ möglich, Bilder zu erzeugen oder einzubinden:
  \begin{description}
  \item[Einbinden:] $\leadsto$ \textsc{Matlab} (später),
  \item[Zeichnen:] $\leadsto$ PGF/TikZ.
  \end{description}

  \vspace*{1em}

  Innerhalb des Textes / einer Formel:
  \tikz[baseline=1pt,scale=.7]{%
  \draw[->] (-.1,0)--(.4,0);
  \draw[->] (0,-.1)--(0,.4);
  \fill[green]
    (0,0)--(.3,.3)--(.3,0)--cycle;
  \draw[thick] (0,0)--(.3,.3)--(.3,0)--cycle;
  \draw (-.1,-.1)--(.4,.4);
  } oder
  \[
    \int_0^1 x\,dx =
    \tikz[baseline=1pt,scale=.7]{%
    \draw[->] (-.1,0)--(.4,0);
    \draw[->] (0,-.1)--(0,.4);
    \fill[green]
      (0,0)--(.3,.3)--(.3,0)--cycle;
    \draw[thick] (0,0)--(.3,.3)--(.3,0)--cycle;
    \draw (-.1,-.1)--(.4,.4);
    } = \frac{1}{2}.
  \]

\begin{verbatim}
  \tikz[baseline=1pt,scale=.7]{%
  \draw[->] (-.1,0)--(.4,0);
  \draw[->] (0,-.1)--(0,.4);
  \fill[green]
    (0,0)--(.3,.3)--(.3,0)--cycle;
  \draw[thick] (0,0)--(.3,.3)--(.3,0)--cycle;
  \draw (-.1,-.1)--(.4,.4);
}
\end{verbatim}
  \begin{textblock*}{5cm}(9.2cm,4.7cm)
\only<2>{%
  \tikz[baseline=1pt,scale=1]{%
  \draw[->] (-.1,0)--(.4,0);
  \draw[->] (0,-.1)--(0,.4);
  \fill[green]
    (0,0)--(.3,.3)--(.3,0)--cycle;
  \draw[thick] (0,0)--(.3,.3)--(.3,0)--cycle;
  \draw (-.1,-.1)--(.4,.4);
}}%
\only<3>{%
  \tikz[baseline=1pt,scale=2]{%
  \draw[->] (-.1,0)--(.4,0);
  \draw[->] (0,-.1)--(0,.4);
  \fill[green]
    (0,0)--(.3,.3)--(.3,0)--cycle;
  \draw[thick] (0,0)--(.3,.3)--(.3,0)--cycle;
  \draw (-.1,-.1)--(.4,.4);
}}%
\only<4>{%
  \tikz[baseline=1pt,scale=3]{%
  \draw[->] (-.1,0)--(.4,0);
  \draw[->] (0,-.1)--(0,.4);
  \fill[green]
    (0,0)--(.3,.3)--(.3,0)--cycle;
  \draw[thick] (0,0)--(.3,.3)--(.3,0)--cycle;
  \draw (-.1,-.1)--(.4,.4);
}}%
\only<5>{%
  \tikz[baseline=1pt,scale=4]{%
  \draw[->] (-.1,0)--(.4,0);
  \draw[->] (0,-.1)--(0,.4);
  \fill[green]
    (0,0)--(.3,.3)--(.3,0)--cycle;
  \draw[thick] (0,0)--(.3,.3)--(.3,0)--cycle;
  \draw (-.1,-.1)--(.4,.4);
}}%
\only<6>{%
  \tikz[baseline=1pt,scale=5]{%
  \draw[->] (-.1,0)--(.4,0);
  \draw[->] (0,-.1)--(0,.4);
  \fill[green]
    (0,0)--(.3,.3)--(.3,0)--cycle;
  \draw[thick] (0,0)--(.3,.3)--(.3,0)--cycle;
  \draw (-.1,-.1)--(.4,.4);
}}
  \end{textblock*}

\end{frame}
% ----------------------------end frame--------------------------------
%----------------------------begin frame------------------------------
\begin{frame}[fragile]
  \frametitle{Schach gefällig?}
  
  \begin{verbatim}
  \tikz[baseline=1.9cm,scale=0.4]{
    \foreach \x in {0,...,7} \foreach \y in {0,...,7}
    {
      \pgfmathparse{int(mod(\x+\y,2))}
      \ifnum\pgfmathresult=0{\draw[fill] (\x,\y)
        rectangle (\x+1,\y+1);}
      \else{\draw (\x,\y)
        rectangle (\x+1,\y+1);}
      \fi
    }
  }
  \end{verbatim}
  \vspace*{-4.5em}
  \begin{figure}[htb]
    \begin{center}
    \tikz[baseline=1.9cm,scale=0.4]{
      \foreach \x in {0,...,7} \foreach \y in {0,...,7}
      {
        \pgfmathparse{int(mod(\x+\y,2))}
        \ifnum\pgfmathresult=0{\draw[fill] (\x,\y)
          rectangle (\x+1,\y+1);}
        \else{\draw (\x,\y)
          rectangle (\x+1,\y+1);}
        \fi
      }
    }
    \end{center}
    \caption{\label{Abb:Schach} Ein Schachbrett mittels PGF/TikZ.}
  \end{figure}

\end{frame}
%----------------------------end frame--------------------------------
%----------------------------begin frame------------------------------
\begin{frame}
  \frametitle{Algorithmen}

  \setcounter{algorithm}{0} % sonst haben wir auf jeder Folie einen
                            % neuen Algorithmus :)
  \begin{algorithm}[H]
  \caption{Das klassische Gram-Schmidt-Verfahren%
    \label{alg:CGS}}
  \begin{algorithmic}[1]
    \REQUIRE $\bfA\in\mathbb{R}^{m\times n}$ mit $\Rang(\bfA)=n$.
    \ENSURE $\bfQ\in\mathbb{R}^{m\times n}$ orthonormal,
    $\bfR\in\mathbb{R}^{n\times n}$ Dreiecksmatrix, $\bfA=\bfQ\bfR$.
    \STATE \alert<2>{$[m,n] \leftarrow \textsf{size}(\bfA)$;}
    \STATE \alert<3>{$\bfQ \leftarrow \zeros(m,n)$;}
    \STATE \alert<4>{$\bfR \leftarrow \zeros(n)$;}
    \FOR{$k = 1:n$}
      \STATE \alert<5,9>{$\bfR(1:k-1,k) \leftarrow
        \bfQ(1:m,1:k-1)^{\mathsf{T}}\bfA(1:m,k)$;}
      \COMMENT{$r_{j,k}=\bfq_j^{\mathsf{T}}\bfa_k$}%
              {$r_{j,k}=\bfq_j^{\mathsf{T}}\bfa_k$}%
      \STATE \alert<6,10>{$\bfv \leftarrow \bfA(1:m,k)-
        \bfQ(1:m,1:k-1)\bfR(1:k-1,k)$;}
      \COMMENT{$\bfv\perp\bfq_1,\ldots$}%
              {$r_{j,k}=\bfq_j^{\mathsf{T}}\bfa_k$}%
      \STATE \alert<7,11>{$\bfR(k,k) \leftarrow \textsf{norm}(\bfv)$;}
      \COMMENT{Berechne Norm $r_{k,k}$ des Vektors}%
              {Berechne Norm $r_{k,k}$ des Vektors}%
      \STATE \alert<8,12>{$\bfQ(1:m,k) \leftarrow \bfv/\bfR(k,k)$;}
      \COMMENT{Normalisiere $\bfq_k$}%
              {Berechne Norm $r_{k,k}$ des Vektors}%
    \ENDFOR
  \end{algorithmic}
\end{algorithm}

\end{frame}
%----------------------------end frame--------------------------------
%----------------------------begin frame------------------------------
\begin{frame}
  \frametitle{Automatismen}

  Warum \LaTeX?
  \begin{itemize}
  \item<1-> \LaTeX\ setzt automatisch Zeilen- und Wortumbrüche.
  \item<2-> \LaTeX\ nummeriert alles für mich.
  \item<3-> \LaTeX\ setzt automatisch mein Literaturverzeichnis.
  \item<4-> \only<4>{\LaTeX\ schreibt automatisch meine
      Arbeiten.}\only<5->{\sout{\LaTeX\ schreibt automatisch meine
      Arbeiten.}}
  \end{itemize}

  \vspace*{2em}

\uncover<6->{%
  Literaturverzeichnis erstellen $\leadsto$ BibTeX mit Datei
  \texttt{.bib}. Woher Einträge?}

\end{frame}
%----------------------------end frame--------------------------------
%----------------------------begin frame------------------------------
\begin{frame}[fragile]
  \frametitle{Literaturverzeichnis}

  \texttt{.bib}-Einträge:

  \begin{quotation}
    \url{http://ams.org/mref}.
  \end{quotation}

  Kopierter Text:\vspace*{-.6em}
\begin{semiverbatim}\tiny
  Linear Algebra and its Applications
  Volume 447, 15 April 2014, Pages 119-132
  An augmented analysis of the perturbed two-sided Lanczos
    tridiagonalization process
  Christopher C. Paige, Ivo Panayotov, Jens-Peter M. Zemke
\end{semiverbatim}

  Bei MRef eingeben, Knopf \texttt{BibTeX} drücken:\vspace*{-.6em}
\begin{semiverbatim}\tiny
@article \{MR3200211,
    AUTHOR = \{Paige, Christopher C. and Panayotov, Ivo and Zemke, Jens-Peter
              M.\},
     TITLE = \{An augmented analysis of the perturbed two-sided \{L\}anczos
              tridiagonalization process\},
   JOURNAL = \{Linear Algebra Appl.\},
  FJOURNAL = \{Linear Algebra and its Applications\},
    VOLUME = \{447\},
      YEAR = \{2014\},
     PAGES = \{119--132\},
      ISSN = \{0024-3795\},
   MRCLASS = \{65F15 (15A21 65G50)\},
  MRNUMBER = \{3200211\},
MRREVIEWER = \{Rafikul Alam\},
       DOI = \{10.1016/j.laa.2013.05.009\},
       URL = \{http://dx.doi.org/10.1016/j.laa.2013.05.009\},
\}
\end{semiverbatim}
  \vspace*{-.6em}
  In BibTeX-Datei hineinkopieren, modifizieren, fertig.

\end{frame}
%----------------------------end frame--------------------------------
\subsection{Die Programmierung in \textsc{Matlab}}
%----------------------------begin frame------------------------------
\begin{frame}[fragile]
  \frametitle{Die Programmierumgebung \textsc{Matlab}}

  \textsc{Matlab} bietet schnellen \& interaktiven Weg, um einfache
  Rechnungen der numerischen linearen Algebra durchzuführen.

  \vspace*{1em}
  
  \textsc{Matlab} = \textsc{Mat}rix \textsc{lab}oratory

  \vspace*{1em}

  Hauptdatentyp: Matrix, Beispiel $\bfA\in\mathbb{R}^{3\times3}$:
\begin{verbatim}
  A = [1 2 3;4 5 6;7 8 9]
\end{verbatim}
  (Auf Einheitslänge normierte) Lösung des homogenen
  Gleichungssystems:
\begin{verbatim}
  n = null(A)
\end{verbatim}
  Umskalieren:
\begin{verbatim}
  x = n/n(1)
\end{verbatim}

\end{frame}
%----------------------------end frame--------------------------------
%----------------------------begin frame------------------------------
\begin{frame}
  \frametitle{Die Programmiersprache \textsc{Matlab}}

  \textsc{Matlab} erlaubt Schleifen, Verzweigungen, Funktionen, \ldots:
  \vspace*{-1em}

  \begin{quotation}
    \lstinputlisting[caption=Der QR-Algorithmus nach Francis]%
      {matlab/FrancisAlgorithmus.m}
  \end{quotation}

\end{frame}
%----------------------------end frame--------------------------------
%----------------------------begin frame------------------------------
\begin{frame}
  \frametitle{\textsc{Matlab}: Schnelle Hilfe}

  Implementierte Techniken der linearen Algebra:
  \begin{itemize}
  \item LR-Zerlegung $\leadsto$ \texttt{lu},
  \item QR-Zerlegung $\leadsto$ \texttt{qr},
  \item $\bfA\bfx=\bfb$ und $\min_{\bfx}\|\bfA\bfx-\bfb\|_2$
    $\leadsto$ $\bfA$\textbackslash$\,\bfb$,
  \item Norm $\leadsto$ \texttt{norm},
  \item Eigenwerte $\leadsto$ \texttt{eig},
  \item SVD $\leadsto$ \texttt{svd}.
  \end{itemize}

  \vspace*{1em}

\uncover<2->{%
  Passende Funktion finden:
  \begin{quotation}
    \texttt{lookfor} \textsf{<Suchbegriff>}
  \end{quotation}}

  \vspace*{1em}

\uncover<3->{%
  Hilfe zu Programmen mittels
  \begin{quotation}
    \texttt{help} \textsf{<Begriff>}
  \end{quotation}}

\end{frame}
%----------------------------end frame--------------------------------
%----------------------------begin frame------------------------------
\begin{frame}
  \frametitle{Grcar-Matrix: \texttt{eig} vs. unser Programm}

  \vspace*{-1em}
  \begin{quotation}
    \lstinputlisting[caption=Eigenwertberechnung und graphische
      Darstellung]{matlab/EigenwertBild.m}
  \end{quotation}

\end{frame}
%----------------------------end frame--------------------------------
%----------------------------begin frame------------------------------
\begin{frame}
  \frametitle{\textsc{Matlab} kann auch schwarz \ldots}

  \begin{center}
    \includegraphics%
      [width=\textwidth,height=.8\textheight,keepaspectratio]%
      {pics/Eigenwerte}
  \end{center}

\end{frame}
%----------------------------end frame--------------------------------
%----------------------------begin frame------------------------------
\begin{frame}
  \frametitle{Von Pseudo-Code zu \textsc{Matlab}}

  Natürlich können wir auch den gezeigten Pseudo-Code in
  \textsc{Matlab} implementieren, hier ist das Ergebnis:
  \begin{quotation}
    \lstinputlisting[caption=Implementation des klassischen
    Gram-Schmidt'schen Orthonormalisierungsverfahrens]%
    {matlab/KlassischesGramSchmidtVerfahren.m}
  \end{quotation}

\end{frame}
%----------------------------end frame--------------------------------
%----------------------------begin frame------------------------------
\begin{frame}[fragile]
  \frametitle{Verwendung des Programms nebst Ausgabe}

  {\footnotesize
  \begin{semiverbatim}
  >> \alert<1>{A = [2 2;2 1;1 0]}
  A =
       2     2
       2     1
       1     0
\uncover<2->{  >> \alert<2>{[Q,R] = KlassischesGramSchmidtVerfahren(A)}
  Q =
       6.66666666666666\alert<5>{6}e-01     6.666666666666667e-01
       6.66666666666666\alert<5>{6}e-01    -3.333333333333333e-01
       3.333333333333333e-01    -6.66666666666666\alert<5>{6}e-01
  R =
       3     2
       0     1}
\uncover<3->{  >> \alert<3>{Q*R}
  ans =
       2     2
       2     1
       1     0}
\uncover<4->{  >> \alert<4>{Q'*Q}
  ans =
       1.000000000000000e+00     \alert<5>{1.110223024625157e-16}
       \alert<5>{1.110223024625157e-16}     1.000000000000000e+00}
  \end{semiverbatim}}

\end{frame}
%----------------------------end frame--------------------------------
\section{Aufbau der Ausarbeitung}
%----------------------------begin frame------------------------------
\begin{frame}
  \frametitle{Grober Aufbau}

  \begin{itemize}
  \item Erster Abschnitt:
  \end{itemize}
  \begin{quote}
  \begin{itemize}
  \item in das Thema adäquat einführen ($\leadsto$ Einführung),
  \item die behandelten Themen motivieren ($\leadsto$ Motivation),
  \item die verwendete Notation erläutern ($\leadsto$ Notation).
  \end{itemize}
  \end{quote}
  \begin{itemize}
  \item Folgende/r Abschnitt/e: wesentlichen Gedanken.
  \item Gesonderter Abschnitt: numerische Beispiele.
  \item Abschluss: Fazit nebst Ausblick.
  \end{itemize}

\end{frame}
%----------------------------end frame--------------------------------
%----------------------------begin frame------------------------------
\begin{frame}
  \frametitle{Inhaltsverzeichnis}

  Ein mögliches Inhaltsverzeichnis sieht wie folgt aus:

  \begin{itemize}
  \item Einführung
    \begin{itemize}
    \item Motivation
    \item Notation
    \end{itemize}
  \item Die wesentliche Idee 1
  \item Die wesentliche Idee 2
  \item Numerische Beispiele
  \item Fazit
  \item Literaturverzeichnis
  \end{itemize}

\end{frame}
%----------------------------end frame--------------------------------
%----------------------------begin frame------------------------------
\begin{frame}
  \frametitle{Plagiate vermeiden!}
  
  Klarstellen, ob Behauptung von Ihnen oder Zitat:
  \begin{quotation}\noindent
    \llap{„}Watkins~\cite[Seite~18]{Watkins:2002} behauptet, dass
    viele physikalische Phänomene durch Differentialgleichungen
    modelliert werden können.“
  \end{quotation}

\uncover<2->{%
  \begin{itemize}
  \item Aussage falsch $\leadsto$ Sie auf der sicheren Seite.
  \item Nicht Lorbeeren für Arbeit anderer einheimsen $\leadsto$
    Plagiat.
  \item Aber: In Einführung Ihren Anteil an Ergebnissen klarstellen!
  \end{itemize}}

  \vspace*{1em}

\uncover<3->{%
  Hilfsmittel wie Wikipedia zitieren:
  \begin{quotation}\noindent
    \url{https://de.wikipedia.org/wiki/Zitieren_von_Internetquellen}.
  \end{quotation}
  Zitieren von Internetquellen unbeliebt $\leadsto$ besser
  Lehrbücher/Journale.}

\end{frame}
%----------------------------end frame--------------------------------
\section{Aufbau eines wissenschaftlichen Vortrages }
%----------------------------begin frame------------------------------
\begin{frame}
  \frametitle{Goldene Regeln}

  \begin{itemize}
  \item 5-7 Zeilen pro Folie
  \item 2 Minuten pro Folie einplanen
  \item keine ausgefeilten Sätze, nur Stichworte
  \item nicht alles vorlesen, Zuhörer können meist lesen
  \item ein Drittel für alle, ein Drittel für Wissende, ein Drittel
    für Spezialisten
  \item “abc-xyz-rule” (englisch)
  \end{itemize}


\end{frame}
%----------------------------end frame--------------------------------
\section{Hilfreiche Internet-Quellen}
%----------------------------begin frame------------------------------
\begin{frame}
  \frametitle{\LaTeX \& \textsc{Matlab}}

  \LaTeX, \TeX:
  \begin{description}
  \item[\href{http://www.latex-project.org/}%
             {http://www.latex-project.org/}] -- \LaTeX-Projekt
  \item[\href{https://www.ctan.org/starter}%
             {https://www.ctan.org/starter}] -- Kurzeinführung
  \item[\href{http://www.tug.org/begin.html}%
             {http://www.tug.org/begin.html}] -- längere
             Kurzeinführung
  \item[\href{https://www.dante.de/tex/TeXAnfaenger.html}%
             {https://www.dante.de/tex/TeXAnfaenger.html}] -- Start
             mit \LaTeX
  \item[\href{http://tex.stackexchange.com/}%
             {http://tex.stackexchange.com/}] -- Fragenportal zu \TeX\
             \& \LaTeX
  \end{description}

  \vspace*{1em}

  \textsc{Matlab}
  \begin{description}
  \item[\href{https://de.mathworks.com/products/matlab.html}%
             {https://de.mathworks.com/products/matlab.html}] --
             \textsc{Matlab}-Hauptseite
  \item[\href{https://de.mathworks.com/matlabcentral/fileexchange/}%
             {https://de.mathworks.com/matlabcentral/fileexchange/}]
             -- Quelle hilfreicher Implementationen \& netter Gimmicks
  \item[\href{https://www.tuhh.de/rzt/software/numerik/matlab/matlab-tsh.html}%
             {https://www.tuhh.de/rzt/software/numerik/matlab/matlab-tsh.html}]
             -- TUHH-Seite zur Installation von \textsc{Matlab} für
             Studierende
  \end{description}

\end{frame}
%----------------------------end frame--------------------------------
%----------------------------begin frame------------------------------
\begin{frame}
  \frametitle{Fazit \& Ausblick}

  Fazit:

  \begin{itemize}
  \item Jetzt sollte hier ein sinniges Fazit stehen.
  \item Sie sollten zumindest einen deutlich kürzeren Vortrag halten.
  \end{itemize}

  \vspace*{1em}

\uncover<2->{%
  Ausblick:

  \begin{itemize}
  \item Wir werden hoffentlich viele schöne Vorträge sehen und einiges
    lernen, am meisten über (und für) uns selbst.
  \end{itemize}}

\end{frame}
%----------------------------end frame--------------------------------
\appendix
\backupbegin
%----------------------------begin frame------------------------------
\begin{frame}
  \frametitle{Appendix}

  Im Appendix können Sie
  \begin{itemize}
  \item Formeln,
  \item Beweise,
  \item Pseudo-Code,
  \item \ldots
  \end{itemize}
  sammeln, von denen Sie denken, dass dazu eine Frage auftauchen
  könnte.

  \vspace*{1em}

  (Dies wird in der Seitenzählung nicht angezeigt $\leadsto$ keiner
  erschrickt sich.)

  \vspace*{1em}

  Ich verwende den Appendix strategisch meist für die folgenden zwei
  Dinge:

\end{frame}
%----------------------------end frame--------------------------------
%----------------------------begin frame------------------------------
\begin{frame}
  \frametitle{Dankeschön!}

\selectlanguage{ngerman}

  \Large
  \begin{center}
    Vielen Dank für die Aufmerksamkeit!

    \vspace*{3em}
    
    \uncover<2->{(Möge Ihr Vortrag ein guter sein.)}
  \end{center}

\end{frame}
%----------------------------end frame--------------------------------
\begin{frame}
  \frametitle{Literatur}

  \bibliography{Proseminar}
  \bibliographystyle{apalike}

\end{frame}

\backupend

\end{document}

%%% Local Variables: 
%%% mode: latex
%%% TeX-master: t
%%% End:
